% ======================================================================
\section{Fourier transforms}\label{app:fourier}
% ======================================================================

With the conventions~\eqref{eq:FTconventions} the transverse integrals that are needed to transform the wave functions and the coefficients from momentum to coordinate space are the following. The Fourier transforms that are needed are, for real $\xi>0$ and real $a,b$,
\begin{align}
&\int\frac{d^2k}{2\pi}\,\frac{k^i}{\v k^2}\,e^{i\v k\cdot\v X}=\frac{iX^i}{\v X^2}\,,\label{eq:FT_K}\\
&\int\frac{d^2k}{2\pi}\frac{d^2p}{2\pi}\,\frac{k^ip^j}{\v k^2\,(a\v k^2+b\v p^2)}\,e^{i\v k\cdot\v X+i\v p\cdot\v Z}=-\frac{X^iZ^j}{\v Z^2\,(b\v X^2+a\v Z^2)}\,,\label{eq:FT_i2}\\
&\int\frac{d^2k}{2\pi}\frac{d^2p}{2\pi}\,\frac{1}{\xi\v k^2-\v p^2}\,e^{i\v k\cdot\v X+i\v p\cdot\v Z}=\frac{1}{\xi\v Z^2-\v X^2}\,,\label{eq:FT_i3}\\
&\int\frac{d^2k}{2\pi}\frac{d^2p}{2\pi}\,\frac{k^m-p^m}{(\v k-\v p)^2}\,\frac{p^l-\xi k^l}{\v p^2-\xi\v k^2}\,e^{i\v k\cdot\v X+i\v p\cdot\v Z}=-\frac{X^m+\xi Z^m}{\v X^2-\xi\v Z^2}\,\frac{X^l+Z^l}{(\v X+\v Z)^2}\,,\label{eq:FT_i4}\\
&\int\frac{d^2k}{2\pi}\frac{d^2p}{2\pi}\,\frac{k^m-p^m}{\v p^2-\xi\v k^2}\,\frac{p^l-\xi k^l}{(\v p-\xi\v k)^2}\,e^{i\v k\cdot\v X+i\v p\cdot\v Z}=-\frac{X^l+Z^l}{\v X^2-\xi\v Z^2}\,\frac{X^m+\xi Z^m}{(\v X+\xi\v Z)^2}\,,\label{eq:FT_i5}\\
&\int\frac{d^2k}{2\pi}\,\frac{k^i}{\v k^2}\,\log\frac{\v k^2}{\mu^2}\,e^{i\v k\cdot\v X}=\frac{iX^i}{\v X^2}\Big[-2\gamma_E-\log\frac{\v X^2\mu^2}{4}\Big]\,,\label{eq:FT_log}
\end{align}
with principal values where a denominator vanishes. Eq.~\eqref{eq:FT_K} gives the Weizs\"acker--Williams field, Eq.~\eqref{eq:kernels}. It follows from $\int_0^{2\pi}d\theta\,\cos\theta\,e^{ikX\cos\theta}=2\pi iJ_1(kX)$ and $\int_0^\infty dt\,J_1(t)=1$, and Eq.~\eqref{eq:FT_log} from the same angular integral and $\int_0^\infty dt\,J_1(t)\log t=\log2-\gamma_E$. The latter relates the $\MS$ logarithm in momentum space to the combination
\begin{align}\label{eq:c0}
2\gamma_E+\log\frac{\v X^2\mu^2_{\MS}}{4}=\log\frac{\v X^2\mu^2_{\MS}}{c_0^2}\,,\qquad c_0=2e^{-\gamma_E}\,,
\end{align}
that appears in coordinate space. Eq.~\eqref{eq:FT_i2} follows from Eq.~\eqref{eq:FT_K} and from the transform of $1/(a\v k^2+b\v p^2)$ in four dimensions. Eq.~\eqref{eq:FT_i3} follows from $\PV\frac1x=\int_0^\infty dt\,\sin(tx)$: the integrals over $\v k$ and $\v p$ are then Gaussian, and the remaining integral over $t$ is of the same form.

Eqs.~\eqref{eq:FT_i3}--\eqref{eq:FT_i5} are the transforms used in~\cite{Lublinsky:2016meo,Radhakrishnan:2026soft}. We have not derived Eqs.~\eqref{eq:FT_i4} and~\eqref{eq:FT_i5} analytically, but we have verified the three principal-value transforms numerically as identities between distributions (script \texttt{scripts/check\_fourier\_pv.py}). Both sides were integrated with a Gaussian test function centred at a generic point, $\exp[-(\v X-\v X_0)^2/2a-(\v Z-\v Z_0)^2/2b]$, and its Fourier transform, respectively. With the representation of the principal value given above and $1/\v q^2=\int_0^\infty ds\,e^{-s\v q^2}$ the transverse integrals are Gaussian on either side, and the remaining integrals over $s$ and $t$ were done numerically. For $\xi=0.37$ and $\xi=2.6$, i.e.\ on both sides of $\xi=1$, and for all components $(m,l)$, the two sides agree to a relative accuracy of $10^{-15}$, while an incorrect right-hand side deviates at order one.
